%Abstract%
\section*{\centering{ABSTRACT}}
\noindent The processing and auditing of Digital Fiscal Bookkeeping for ICMS
and IPI (EFD ICMS IPI) documents face operational bottlenecks in both the
public and private sectors. The hierarchical plain text structure of these
files, the high volatility of legal SPED layouts, and the absence of an
official machine-readable specification render traditional hard-coded
systems inflexible and error-prone. To solve this problem, this study
develops a dynamic, metadata-driven ingestion pipeline. The proposed
architecture extracts business rules from tax documentation and structures
them into JSON configuration files, decoupling data processing from the source
code. The solution consists of an Orchestrator that distributes the workload
into batches and manages the parallel execution of in-memory dynamic workers,
which parse the textual hierarchy and perform structured persistence into a
\textit{PostgreSQL} relational database. Empirical validation demonstrated
zero information loss and preserved referential integrity through reverse
reconstruction tests of the original data. The solution also showed high
resilience in handling corrupted files and ensured database transaction
atomicity. It is concluded that the metadata-driven approach provides high
flexibility regarding changes in tax legislation, while the identified
limitations offer substantial room for future development and tool expansions.

\vspace{\baselineskip}

\noindent {\bf Keywords:} EFD ICMS IPI.
SPED. Data Engineering.
Data Ingestion.
Metadata-Driven Architecture.
Relational Database.